
\begin{center}
\Huge
	Opgaveregning
\end{center}

\subsection*{Opgave 1}
Bestem koordinatsættet til toppunktet for følgende andengradspolynomier

\begin{align*}
	&1) \ p(x) = x^2 + 2x + 4 &  &2) \ p(x) = -2x^2 + 6x - 10 \\
	&3) \ p(x) = 3x^2 - 5x + 2 &  &4) \ p(x) = \frac{1}{2}x^2 + x - 2 \\
\end{align*}

\subsection*{Opgave 2}

Et andengradspolynomium $p$ er givet ved
\begin{align*}
	p(x) = 2x^2 -4x + c
\end{align*}
\begin{enumerate}[label=\roman*)]
	\item Bestem tallet $c$, så $p$ har toppunkt i $(1,32)$.
	\item For hvilke værdier af $c$ ligger toppunktet under $x$-aksen?
\end{enumerate}

\subsection*{Opgave 3}
To ligninger er givet ved
\begin{align*}
	2x + 4y &= 10\\
	x + 6y &= -6\\
\end{align*}
\begin{enumerate}[label=\roman*)]
	\item Bestem løsningen til ligningssystemet.
\end{enumerate}

\subsection*{Opgave 4}
Løs følgende ligninger ved brug af nulreglen
\begin{align*}
	&1) \ 10x^2 - 100x = 0 &  &2) \ e^{6x^2 + 24x} = 1
\end{align*}

\subsection*{Opgave 5}
Bestem graden af følgende polynomier

\begin{align*}
	&1) \ p(x) = x^6+ 2x^3 + 4x + 1 & &2) \ p(x) = x^5+ 2x^5\\
	&3) \ p(x) = x^2x^7+x^4x^1 & &4) \ p(x) = 4(x^2)^5
\end{align*}

\subsection*{Opgave 6}
Bestem rødderne for følgende polynomier ved at faktorisere

\begin{align*}
	&1) \ p(x) = x^2 - 2x - 8 & &2) \ p(x) = x^2 +7x + 10  \\
	&3) \ p(x) = x^2 - 6x + 5  & &4) \ p(x) = x^2 - 3x - 4 \\
	&5) \ p(x) = x^2 + x - 2 & &6) \ p(x) = x^2 + 4x + 3 \\
\end{align*}